\PassOptionsToPackage{table}{xcolor}
\documentclass{article}
\usepackage{iclr2027_conference,times}
\usepackage{hyperref}
\usepackage{url}
\usepackage{graphicx}
\usepackage{booktabs}
\usepackage{amsmath}
\usepackage{xcolor}
\definecolor{linkblue}{HTML}{1F5A94}
\hypersetup{
  colorlinks=true,
  linkcolor=linkblue,
  citecolor=linkblue,
  urlcolor=linkblue
}
\usepackage{amsfonts}
\usepackage{amssymb}
\usepackage{amsthm}
\usepackage{enumitem}
\usepackage{pifont}
\DeclareMathSizes{9}{10}{7}{5}

% Theorem environments and shared math notation.
\newtheorem{theorem}{Theorem}
\newtheorem{assumption}{Assumption}
\newtheorem{proposition}{Proposition}
\theoremstyle{definition}
\newtheorem{definition}{Definition}
\theoremstyle{plain}
\newcommand{\E}{\mathbb{E}}
\newcommand{\Var}{\mathrm{Var}}
\newcommand{\norm}[1]{\lVert #1 \rVert}

% Tighten float/caption spacing to fit the 9-page limit (formatting only).
\setlength{\textfloatsep}{8pt plus 2pt minus 2pt}
\setlength{\floatsep}{8pt plus 2pt minus 2pt}
\setlength{\intextsep}{8pt plus 2pt minus 2pt}
\setlength{\abovecaptionskip}{4pt}
\setlength{\belowcaptionskip}{2pt}

% Compact captions: keep plots readable while avoiding body-sized caption blocks.
\makeatletter
\long\def\@makecaption#1#2{%
  \vskip\abovecaptionskip
  {\footnotesize
    \sbox\@tempboxa{\textbf{#1:} #2}%
    \ifdim \wd\@tempboxa >\hsize
      \textbf{#1:} #2\par
    \else
      \global\@minipagefalse
      \hb@xt@\hsize{\hfil\box\@tempboxa\hfil}%
    \fi}
  \vskip\belowcaptionskip}
\makeatother

% -----------------------------------------------------------------------
% Working title — change before submission
% -----------------------------------------------------------------------
\title{Weights Have Traits: Estimating Fine-Tuning Lineage in LLMs}

% Authors hidden for double-blind review.
% \iclrfinalcopy  % Uncomment for camera-ready only.

\newcommand{\fix}{\marginpar{FIX}}
\newcommand{\new}{\marginpar{NEW}}

% Placeholder macro for unwritten sections — shows in red, visible in draft
\newcommand{\TODO}[1]{\textcolor{red}{\textbf{[TODO: #1]}}}

% Shorthand for figure path prefix
\newcommand{\figpath}{figures}

\begin{document}
\maketitle

% -----------------------------------------------------------------------
\begin{abstract}
While over a million deep learning models are hosted online, most are fine-tuned descendants from a small set of ancestors.  The lineage of these models is only sometimes documented fully, even though provenance, auditing, and safe deployment all depend on proper documentation. In response to this issue, we hypothesize that fine-tuning acts as an evolutionary process that leaves signal in weight space, and show that a model \textit{phylogeny} (evolutionary tree) can be estimated from model weights alone and scored against a known ground-truth phylogeny, recovering all true clades (branches) in 85--100\% of runs.  We treat individual weight layers as gene-like estimators of lineage, revealing where evolutionary signal is concentrated and allowing each layer to produce its own phylogeny.  We are able to show that we can accurately estimate the phylogeny using only a percent of weights (1--3\% in our experiments), suggesting that lineage can be recovered from a small, architecture-specific weight subset. Moreover, we show that this phylogeny estimation is associated with known mathematical concepts from phylogenetics: the Atteson margin and four-point additivity. Beyond phylogeny estimation, we report two additional empirical findings. First, weight-space proximity predicts behavioral similarity on several held-out tasks involving truthfulness, knowledge, reasoning, and common sense, though the strength and statistical support depend on the probe. Second, weight-based phylogeny estimation can match or outperform behavior-based estimation, particularly after low-rank updates. We finally compare our weight-based method to a state-of-the-art behavior-based method on both our controlled experiments and a real, albeit small, modern model phylogeny. On the matched controlled trees, weights recover 100\% of clades versus PhyloLM's 35\% under rank-8 LoRA and 100\% versus 52\% under full fine-tuning, while both methods recover all three documented splits in the real Mistral family.

\end{abstract}

% -----------------------------------------------------------------------
\section{Introduction}
\label{sec:intro}




While there are over one million models available on HuggingFace for deep learning and a handful of frontier
labs constantly releasing updated versions, most models are
fine-tuned descendants of a few ancestors.  Very few of those models transparently document their full model lineage \citep{horwitz2025chartatlasworldsmodels}. Benchmarks can reveal how a model performs on a task, but not how a model originated, how it differs from its sibling models, or how it may
fail. In contrast, training lineage can reveal much about model provenance, auditing, and safety~\citep{behavioral2026provenance, nikolic2026model, tensorguard2025}.

We, like other prior works~\citep{horwitz2025mother, yu2025neural, wu2026llmdna,gallagher2026}, observe that fine-tuning is a path-dependent, branching process. Each
child model inherits its parent's  updated weights, nudged towards a certain loss and/or data set. This is
precisely an evolutionary process, as described in \cite{yax2025phylolm}, and moreover, the evolutionary path leaves signal in weight space. We treat a line or genealogy of fine-tuned models
as a phylogeny (i.e. an evolutionary tree~\cite{saitou1987neighbor}) and ask whether this phylogeny can be \emph{estimated} from the model weights alone.

As motivation, consider a firm that forks a base model into finance, IT, and fraud variants, each fine-tuned until
dozens of versions circulate with no recorded lineage (Fig.~\ref{fig:overview}A). A practitioner needs
to know which to build on, where two versions diverged, and where each fails. Phylogenetic algorithms
were built for exactly this inference in biology~\citep{saitou1987neighbor, robinson1981comparison},
and we adapt them to neural net model weights. We estimate a phylogeny from model weights and score it against a \emph{known}
ground-truth phylogeny (Fig.~\ref{fig:overview}A), link the changes in weights to changes in behavior (Fig.~\ref{fig:overview}B), and ask whether diagnostic statistics indicate when phylogeny estimation will succeed even without access to the true tree (Fig.~\ref{fig:overview}C).

\begin{figure}[t]
  \begin{center}
    \includegraphics[width=\linewidth]{\figpath/fig1_paper_overview.pdf}
  \end{center}
  \caption{Overview of our three questions: Can model weights recover known training relationships? Do nearby models behave alike? Do tree-like distances signal reliable recovery without knowing the true tree?}
  \label{fig:overview}
\end{figure}



Our primary contribution is that analyzing weight changes between fine-tunings allows accurate phylogeny estimation with latent ancestors. We use a fitted Atteson diagnostic and label-free four-point additivity as empirical diagnostics to characterize when phylogeny estimation succeeds. As a result of these weight-based estimates, we can further show two behavior-related consequences. In summary we show:
\begin{itemize}[leftmargin=1.3em,itemsep=1pt,topsep=2pt,parsep=0pt]
  \item \textbf{Phylogeny estimation from leaf model weights against ground truth.} Guided by the geometry of
    fine-tuning updates, our distance metric and summary-statistic choices estimate the model phylogeny from
    weights alone. Our controlled experiments recover \textit{all} true clades (branches) in $85$--$100\%$ of runs
    ($92$--$100\%$ clade recovery). We evaluate against known ground-truth phylogenies and show that a subset containing about 1--3\% of the scalar weights can accurately estimate the known phylogenies. Treating weight tensors as gene-like estimators also yields a separate phylogeny from each tensor and highlights where lineage signal is strongest. We show that our weight-based method accurately recovers all the splits in a small but real Mistral phylogeny.
  \item \textbf{Scenarios where weight-based estimation outperforms behavior-based estimation.}
    On the same 46 controlled trees, weights recover 100\% of clades versus 35\% for
PhyloLM~\citep{yax2025phylolm} under rank-8 LoRA, and 100\% versus 52\% under full fine-tuning.
    On a documented Mistral family with large RLHF/DPO forks, both methods
    recover all three true splits. Thus behavior-only estimation can succeed after behaviorally large
    updates, while weight access remains substantially more informative for smaller updates and
    requires no output sampling.
  \item \textbf{Evidence of similar weights signaling similar behavior.} Changes in weights can be used to predict similarity in model behavior. In particular, we show that models that are closer in weight-space behave more similarly in behavior-space on \textit{held-out} tasks such as truthfulness, reasoning, knowledge, and common sense.
\end{itemize}
% Submission TODO: add the anonymous code-and-artifact URL here once the blind snapshot exists.
Our main controlled topology results cover Flan-T5-base (250M) and Llama-3.2-1B, with both full
fine-tuning and variants of LoRA.






% -----------------------------------------------------------------------
% Swati Rallapalli's related-work draft; keep this as a standalone source.
\input{new_related}

% -----------------------------------------------------------------------x
\section{Framework}
\label{sec:framework}
% -----------------------------------------------------------------------

We first define our terms, notation, and statistical framework. A \emph{phylogeny} $\mathcal{P}$ is a tree
whose leaves are the observed models and whose internal nodes are latent common ancestors; a
\emph{clade} is the set of leaves descending from one such ancestor (one branch of the tree). We analyze
single-root model families but observe only their terminal models. Thus,  we estimate the unrooted leaf topology with internal ancestors latent. We form a pairwise distance matrix and apply neighbor-joining (NJ)
\citep{saitou1987neighbor} to estimate the phylogeny.

\textbf{Phylogeny estimation.}\label{sec:framework:phylo}
Assume we have a model $g(\theta)$ where $\theta$ is composed of multiple layers $\theta_\ell$. To estimate $\widehat{\mathcal{P}}$: (1) choose a set of layers $\theta_\mathcal{L}$ and a distance metric $d$;
(2) for each layer $\ell$ compute $D_{i,j}^{(\ell)} = d(\theta^{(\ell)}_i, \theta^{(\ell)}_j)$ between models
$i,j$; and (3) apply NJ to each selected layer's matrix to obtain a layer-level tree. We either
analyze these layer-level trees separately or combine them in tree space, as described below. Analyses
that require one pairwise weight distance instead use the arithmetic mean of the active layer-level
distance matrices.

Step (3) lets different layers yield different estimates of $\widehat{\mathcal{P}}$, which we use to
assess each layer's lineage signal. We compare estimated and true phylogenies using Robinson--Foulds (RF) distance \citep{robinson1981comparison}, which counts disagreeing bipartitions as false positives (FP) and false negatives (FN). We restrict our analysis to unrooted phylogenies with at least four leaves, since three leaves have only one possible unrooted topology and therefore no nontrivial internal split to recover. We compare estimated and ground-truth phylogenies using Robinson--Foulds (RF) distance \citep{robinson1981comparison}, which measures disagreement between their splits. Because NJ returns binary trees while our ground-truth phylogenies may contain \textit{polytomies} (nodes with more than two children), NJ must introduce additional splits to resolve them. RF counts a split present in the estimated tree but absent from the ground-truth tree as a false positive (FP); thus, splits introduced solely to resolve a polytomy may be counted as FPs even though they do not contradict the known lineage. We therefore focus on false negatives (FN), which count ground-truth clades missed by the reconstruction. We also report clade recovery, $1-\mathrm{FN}/|\text{true splits}|$, the fraction of ground-truth clades recovered, and polytomy-aware exact recovery (PAER), the fraction of runs in which all ground-truth clades are recovered ($\mathrm{FN}=0$). Separately, we report the standardized within-run branch correlation $r_{\mathrm{branch}}$ between weight distance and a same-branch indicator to assess whether weight distance reflects lineage relationships.



\textbf{Distance metric.}\label{sec:framework:metric}
We flatten each weight tensor to a vector and compare models with vector distances rather than matrix
norms, for portability across architectures, considering cosine, correlation, $\ell_1$, and $\ell_2$
(the last used in MoTHer~\citep{horwitz2025mother}). We adopt cosine distance ($1$ minus cosine
similarity) throughout, for three reasons. It is bounded and scale-invariant across layers, so no single
high-magnitude layer dominates the aggregate, unlike $\ell_1$ or $\ell_2$. It is the prevailing choice in
recent provenance and fingerprinting work \citep{zeng2024huref, behavioral2026provenance,
zhu2025independence}, keeping our distances comparable to that literature. Finally, cosine distance has a simple geometric interpretation: descendant models whose weight vectors remain similarly oriented in parameter space are considered closer. For LoRA-trained trees, we merge the ordered root-to-leaf adapter chain into each checkpoint before comparing weights.



\textbf{Per-layer estimated phylogenies.}
We treat each weight layer as an independent estimator of lineage, analogous to using individual
genes to estimate biological phylogenies~\citep{maddison1997gene}. For a given training tree and
layer, we compute the cosine-distance matrix across its observed models and apply NJ, producing one
estimated phylogeny for that layer. The per-layer profile keeps these estimates separate: it scores
each layer's tree against the known training tree and then, for each aligned layer, averages the
resulting FN rate across the 46 trees retained for topology recovery. It never averages distance
matrices across layers or combines different layers' trees. We exclude layers for which all pairwise
cosine distances are zero, since they provide no lineage information under this metric.

When one aggregate topology is required, we instead compute the Billera--Holmes--Vogtmann (BHV)
Fr\'echet mean of the layer-level trees~\citep{billera2001geometry, willis2019confidence}. This is a
tree-space mean: it returns the point in BHV tree space that minimizes the summed squared distance
to the individual layer-level trees, rather than first averaging their distance matrices. Our
primary analysis uses the original cosine distances; we additionally normalize each layer by its
mean pairwise distance as a sensitivity analysis. For analyses that require a single pairwise
weight-distance matrix---including direct lineage, four-point additivity, the Atteson diagnostic,
and weight--behavior associations---we use the arithmetic mean of the nonzero active-layer distance
matrices. Behavioral similarity is likewise the arithmetic mean of prompt-paired output-embedding
cosine similarities. Thus the per-layer profile, aggregate topology, and pairwise behavioral
analysis use separate, explicit aggregation rules.



\textbf{When estimation succeeds.}\label{sec:framework:theory}
Model lineage can be probed at two levels: pairwise ordering (e.g.\ whether two models sit on the same
branch) or whole-tree recovery, the former being easier. For ordering, we pool all pairwise model
comparisons within each valid run and test with the Mann--Whitney $U$ statistic. This branch-level
diagnostic is simpler than either pairwise parent--child detection~\citep{yu2025neural}, MoTHer's
directed observed-node tree recovery~\citep{horwitz2025mother}, or our latent-ancestor leaf topology. Recovering the full
\textit{phylogeny} is harder but well studied in phylogenetics. The NJ method is exact when the distance matrix is
additive, a condition model-weight distances satisfy only approximately. Atteson's theorem makes the
tolerance precise.
\begin{theorem}[\citealp{atteson1999performance}]\label{thm:atteson}
Let $D_T$ be the additive leaf metric of an edge-weighted binary tree $T$ whose edges have strictly
positive lengths $\{L_s\}$, and let $D$ be an observed dissimilarity map. If
$\lVert D-D_T\rVert_\infty < \tfrac12\min_s L_s$, then neighbor-joining applied to $D$ returns $T$.
\end{theorem}
For the theorem's fixed true metric, the ratio
$m:=\min_s L_s/(2\lVert D-D_T\rVert_\infty)>1$ is a sufficient condition of exact recovery. It is
conservative, so NJ can recover the topology when $m\leq1$. In our controlled experiments the true topology is
known but its edge lengths are not. We fit nonnegative edge lengths on that topology to obtain
$D_{\mathrm{fit}}$ and report the analogous diagnostic
$\widetilde m=\min_s\widehat L_s/(2\lVert D-D_{\mathrm{fit}}\rVert_\infty)$. Because
$D_{\mathrm{fit}}$ is estimated from $D$ and some truth trees contain polytomies, $\widetilde m$ is an approximation rather than the true Atteson condition.  Because of this fitting approximation, we also
adopt a label-free proxy computed from $D$ alone, based on the classical four-point condition for tree
metrics~\citep{sempleSteel2003phylogenetics}.

\textit{Four-point additivity} ($A$). Non-additivity of a distance matrix is measured by the quartet
$\delta$-score of \citet{holland2002delta}. For each quartet of leaves the three pairwise-distance sums,
sorted $s_1\ge s_2\ge s_3$, give
\[
  \delta_q=\frac{s_1-s_2}{s_1-s_3}\in[0,1] ,
\]
which is $0$ exactly when the four-point condition $s_1=s_2$ holds \citep{sempleSteel2003phylogenetics}
(the additive, tree-metric case) and $1$ at the maximally non-treelike quartet (the two \emph{smallest}
sums equal, the largest strictly greater). We report its complement, averaged over the $Q$ quartets,
\[
  A=\frac{1}{Q}\sum_{q=1}^{Q}\bigl(1-\delta_q\bigr)
   =\frac{1}{Q}\sum_{q=1}^{Q}\frac{s_2^{(q)}-s_3^{(q)}}{s_1^{(q)}-s_3^{(q)}}\in[0,1] ,
\]
so that $A=1$ is perfectly additive and larger $A$ is associated with more accurate phylogeny estimation. Because $A$ requires no labels,
it serves as our proxy when the true tree is unknown. For evaluation against a known tree we
additionally report $A$ over only the \emph{informative} quartets (those the truth resolves into two
cherries), since quartets that fall inside a polytomy carry no split to recover and only dilute the
all-quartet mean. We relate $A$ and the fitted Atteson diagnostic to phylogeny estimation accuracy in \S\ref{sec:results}.


% -----------------------------------------------------------------------
\section{Experiments}
\label{sec:experiments}
% -----------------------------------------------------------------------
We structure our experiments around three questions: whether the model lineage or phylogeny can be
estimated from weight changes; whether those changes predict downstream held-out behavior; and
how behavior-based phylogeny differs from weight-based.

\paragraph{Controlled and real training trees.}
\label{sec:exp:trees}
% Base: Flan-T5-base (250M). Bernoulli branching, 4 tasks, 50 replications,
% 5--8 leaves, known topology by construction.
% 2000 steps, lr=3e-4, batch=8, grad_accum=4, bf16.
% Controlled design enables exact evaluation against the latent-ancestor leaf topology.

To create a controlled experiment, we generate stochastic model training trees beginning from a root
model (process detailed in App.~\ref{app:details}). The topology-recovery cohort contains the 46 trees
with a nontrivial unrooted truth split, each with 4--10 observed leaves; the behavioral analysis uses
all 50 generated trees. We randomly assign a task/dataset pair to each edge, and at each split
fine-tune a child model from its parent on that pair. LoRA chains are sequentially merged so the
primary LoRA leaves are standalone full-weight checkpoints; the separately marked all-projections
ablation is analyzed in cumulative update space. We then estimate a phylogeny from the \textit{model
leaves} and compare it to the known training tree, testing both direct lineage and phylogeny
estimation.

We vary the root model, task types, and trained layers. Our primary evaluation uses an
encoder--decoder model (Flan-T5-base, 250M) and a decoder-only model (Llama-3.2-1B), fine-tuning on
Q\&A, classification, translation, and summarization. We also run auxiliary analyses on Pythia-160M and Llama-3.1-8B. For Flan-T5, we run full fine-tuning and four LoRA \citep{hu2022lora}
scopes (key only; query/key/value; query/key/value/output; all attention and feed-forward projections).
For Llama-3.2-1B, we run full fine-tuning and query/key/value LoRA at ranks 8 and 64. For each base
model, we use the same task instructions. We also
validate our method on three documented model lineages from HuggingFace. We highlight one in the
main text and report the other two in App.~\ref{app:organic-additional}.

\paragraph{Direct lineage.}
MoTHer~\citep{horwitz2025mother} assembles pairwise relatedness and direction scores into a directed
model tree, while Neural Phylogeny~\citep{yu2025neural} identifies parent--child relationships. Our
auxiliary ordering analysis instead works at the branch level among observed leaves. Within each
training tree, two leaves are labeled ``same branch'' when their first ancestor below the root is
the same; all other leaf pairs are labeled ``cross branch.'' We test whether cross-branch distances
are larger with a one-sided Mann--Whitney $U$ test and report the corresponding rank-biserial effect
$r_{\mathrm{rb}}=2U/(n_{\mathrm{cross}}n_{\mathrm{same}})-1$. Thus $r_{\mathrm{rb}}=1$ means that
every cross-branch pair is farther apart than every same-branch pair, whereas zero indicates no
rank separation. We also report the point-biserial correlation
$r_{\mathrm{branch}}=\operatorname{Cor}(\mathbf{1}\{\mathrm{same\ branch}\},d)$; stronger lineage
signal is negative because same-branch pairs are expected to be closer. Both statistics are first
computed separately within each of the 26 trees that contain both pair classes, so trees with more
leaves do not dominate through having more pairs. Across trees, we report the arithmetic mean and
sample SE of $r_{\mathrm{rb}}$, and Fisher $z$-average $r_{\mathrm{branch}}$\footnote{For each run we
compute $z=\operatorname{arctanh}(r)$, average on the $z$ scale, and back-transform with
$r=\tanh(\bar z)$.}.

\paragraph{Phylogeny estimation.}
\label{sec:exp:weights}
Besides direct lineage, we test whether our method estimates the phylogeny. From \S\ref{sec:framework:theory}, accurate estimation should be possible when the distance matrix is sufficiently additive, for which the label-free four-point additivity $A$ is a practical, label-free proxy. For each tree we report RF, FN, clade recovery, and polytomy-aware exact recovery (PAER; $\mathrm{FN}=0$), relating them to the fitted Atteson diagnostic $\widetilde m$ and $A$ in \S\ref{sec:results}. The per-layer trees also let us identify which layers carry the most inherited signal, via the highest- and lowest-FN layers per architecture.


\paragraph{Predicting held-out behavior.}
\label{sec:exp:behavioral}
We next estimate how strongly weight distance predicts behavioral similarity on probe tasks using the trained, controlled experiment trees. For the translation probe, we obtain one deterministic greedy response per prompt and model, embed each output with a small sentence-transformer (\texttt{all-MiniLM-L6-v2}, $384$-d), and average prompt-paired cosine similarity. We also evaluate HellaSwag, ARC-C, MMLU, and TruthfulQA on the same trees using the corresponding prompt-paired semantic endpoint. We then report the random-effects DerSimonian--Laird \citep{dersimonianlaird1986} pooled correlation $r_{\mathrm{DL}}$ between weight distance and output-embedding similarity, with both quantities $z$-scored within run (placing all models and architectures on one scale). Other generation details are reported in App.~\ref{app:details}.




\paragraph{Comparison to behavior-only estimation.}
\label{sec:exp:blackbox}
Sometimes weights are unavailable and provenance must be inferred from outputs alone, like the method PhyloLM \citep{wu2026llmdna,yax2025phylolm}.  Because our documented HuggingFace family (\S\ref{sec:exp:trees}) has a known topology, we can score behavior-only estimation against its documented reference lineage. On that family we run the behavior-based PhyloLM method~\citep{yax2025phylolm}, using its released genes and genetic distance on the same checkpoints, and compare it against our weight-based estimation on the identical documented lineage (details in App.~\ref{app:blackbox}). Our controlled comparison runs PhyloLM on every tree for the Llama rank-8 LoRA, rank-64 LoRA, and full-FT conditions (50 reconstructions per condition). We compare on trees with at least one nontrivial unrooted truth split. The encoder--decoder Flan groups are outside PhyloLM's decoder-only scope.


% -----------------------------------------------------------------------
\section{Results}
\label{sec:results}
% -----------------------------------------------------------------------

%\subsection{Weights contain both direct ancestry and phylogenetic signal}

We examine the results of our controlled and real data experiments with our encoder-decoder and decoder models, LoRA trainings, and different task fine-tunings and data sets.


\paragraph{Direct lineage.}
We first determine how well direct lineage estimation (e.g., parent--child or within-branch) performs under different conditions. There are $n_{ord}=26$ eligible trees for this analysis that contain at least one pair of leaves from the same root branch and at least one pair from different root branches. Table~\ref{tab:variants} reports the Mann--Whitney rank-biserial effect and the within-run branch correlation $r_{\mathrm{branch}}$ between cosine distance and a same-branch indicator. Every mean rank-biserial effect is positive and every $r_{\mathrm{branch}}$ is negative. This implies that ordering in a tree has a strong signal in the weights in our experiments.

\begin{table}[h]
\caption{Lineage and recovery by training regime. Topology recovery uses the
  BHV Fr\'echet mean of per-layer NJ trees; pairwise lineage statistics use the mean active-layer
  distance matrix. LoRA uses sequentially merged full weights and only layers with nonzero
  between-model cosine distance;
  $\dagger$ marks the cumulative update-space all-projections experiment.
  PAER means $\mathrm{FN}=0$. Llama weight and PhyloLM values share the same $n_{\rm rec}=46$
  eligible trees; Flan is outside PhyloLM's scope. $^\star$Ceiling cells correspond to 117/117
  recovered truth splits and 46/46 PAER trees (two-sided exact 95\% binomial intervals
  96.9--100\% and 92.3--100\%, respectively).}
\label{tab:variants}
\begin{center}
\scriptsize
\setlength{\tabcolsep}{1.1pt}
\resizebox{\linewidth}{!}{%
\begin{tabular}{lrrrrrrr}
\toprule
Variant & $n_{\rm rec}/n_{\rm ord}$ & \shortstack{Rank-bis.\\(SE)} & \shortstack{$r_{\mathrm{branch}}$\\(SE)} & \shortstack{Weight clade\\(SE)} & \shortstack{PhyloLM clade\\(SE)} & \shortstack{Weight PAER\\(SE)} & \shortstack{Weight RF (SE);\\FN (SE)} \\
\midrule
\multicolumn{8}{l}{\textbf{Encoder--decoder: Flan-T5-base} ($d=768$)} \\
\multicolumn{8}{l}{\quad\textit{LoRA adapter scope}} \\
\quad LoRA: key only                    & 46/26 & 0.52 (0.09) & $-0.53$ (0.07) & \textbf{92\%} (3\%) & n/a & \textbf{85\%} (5\%) & 2.39 (0.25); 0.17 (0.06) \\
\quad LoRA: q/k/v                       & 46/26 & 0.56 (0.09) & $-0.62$ (0.07) & \textbf{99\%} (1\%) & n/a & \textbf{98\%} (2\%) & 2.09 (0.21); 0.02 (0.02) \\
\quad LoRA: q/k/v/o (all attention)     & 46/26 & 0.56 (0.09) & $-0.62$ (0.07) & \textbf{99\%} (1\%) & n/a & \textbf{98\%} (2\%) & 2.09 (0.21); 0.02 (0.02) \\
\quad LoRA: q/k/v/o + FFN projections$^\dagger$ & 46/26 & 1.00 (0.00) & $-0.91$ (0.02) & \textbf{96\%} (2\%) & n/a & \textbf{89\%} (5\%) & 2.26 (0.21); 0.11 (0.05) \\
\multicolumn{8}{l}{\quad\textit{Full fine-tuning}} \\
\quad Full fine-tuning                  & 46/26 & 0.69 (0.08) & $-0.71$ (0.06) & \textbf{100\%}$^\star$ & n/a & \textbf{100\%}$^\star$ & 2.04 (0.19); 0.00 (0.00) \\
\midrule
\multicolumn{8}{l}{\textbf{Decoder-only: Llama-3.2-1B} ($d=2048$)} \\
\multicolumn{8}{l}{\quad\textit{Controlled exp. trees}} \\
\quad LoRA qkv ($r{=}8$)                 & 46/26 & 0.63 (0.08) & $-0.64$ (0.07) & \textbf{100\%}$^\star$ & 35\% (5\%) & \textbf{100\%}$^\star$ & 2.04 (0.19); 0.00 (0.00) \\
\quad LoRA qkv ($r{=}64$)                & 46/26 & 0.65 (0.08) & $-0.67$ (0.07) & \textbf{100\%}$^\star$ & 24\% (5\%) & \textbf{100\%}$^\star$ & 2.04 (0.19); 0.00 (0.00) \\
\quad Full fine-tuning                   & 46/26 & 0.64 (0.08) & $-0.66$ (0.07) & \textbf{100\%}$^\star$ & 52\% (5\%) & \textbf{100\%}$^\star$ & 1.83 (0.17); 0.00 (0.00) \\
\bottomrule
\end{tabular}
}
\end{center}
\end{table}



\paragraph{Phylogeny estimation.}
\label{sec:results:weighttree}
In addition to ordering, the BHV Fr\'echet means estimate the phylogeny, as shown in Table~\ref{tab:variants}, with respect to clade recovery, polytomy-aware exact recovery (PAER; $\mathrm{FN}=0$), and RF and FN for all eligible phylogenies. In the controlled experiments, clade recovery ranges from 92--100\% and PAER from 85--100\%. The merged LoRA variants recover 92--99\% of clades for Flan and 100\% for Llama; full fine-tuning recovers 100\% and 100\%, respectively. On the same 46 trees with a nontrivial truth split, weights recover 100\% versus PhyloLM's 35\% under rank-8 LoRA, 100\% versus 24\% at rank 64, and 100\% versus 52\% under full fine-tuning.

Treating pairwise layer distance as gene-like estimators reveals that a small attention-weight subset can estimate the phylogeny as well as the all-layer aggregate in the full-fine-tuning regime. On the same 46 matched-topology Flan runs, the encoder self-attention key subset and the all-weight estimate both recover 100\% of clades with 100\% PAER (App.~\ref{app:layers}). This subset contains 7.08M of 247.58M scalar weights (2.86\%, approximately 3\%).  The result is even more striking for Llama with 100\% of clades recovered and 100\% with LoRA training at rank 8. The individual-layer profile in App.~\ref{app:layertrace} shows that the attention layers are strong carriers of lineage signal, although this is architecture dependent.

Finally, Figure~\ref{fig:coherence_recovery} shows simulation-level, layer-subset, and individual-layer phylogeny estimation summaries from the same full-weight trees as Table~\ref{tab:variants}, together with the matched PhyloLM estimates for the three Llama conditions. Weight-based estimation of clade recovery ranges from 92--100\%, including 100\% for Flan full fine-tuning and 100\% for Llama full fine-tuning. Layer-subset estimates recover 90--100\% of true clades. The matched PhyloLM estimates recover 24--52\% of clades.  This indicates that there are scenarios where weight-based estimation of lineage is more accurate than behavior based estimation.

\begin{figure}[t]
  \begin{center}
    \includegraphics[width=\linewidth]{\figpath/fig4_coherence_atteson.pdf}
  \end{center}
  \caption{Mean true-clade recovery against
    \textbf{(A)} the fitted Atteson diagnostic $\widetilde m$ on a logarithmic x-axis and
    \textbf{(B)} label-free quartet
    additivity $A$. Large filled symbols summarize seven Flan and Llama simulations: full fine-tuning
    plus the five complete merged-weight LoRA conditions.
    Every mark is an average over the same 46 topology-eligible trees used in
    Table~\ref{tab:variants}; vertical intervals on simulation, PhyloLM, and
    subset summaries are $\pm1$ SE. Color identifies the seven primary training conditions only;
    auxiliary layer subsets and individual $q/k/v$ matrices are shown in gray. Dark dashed curves are pooled fractional-logit
    fits to all displayed estimator summaries.}
  \label{fig:coherence_recovery}
\end{figure}


\paragraph{Predicting held-out behavior.}
\label{sec:results:behavioral}

While estimating phylogeny from weights is useful for safety and provenance, it can be of even more use if similarity in weights signals similarity in behavior. Inspired by \citet{cloud2025subliminal}'s subliminal-learning experiment, we test
whether models closer in weight space produce more similar outputs on \textit{held-out} tasks, using translation
(on a new set of trees that never trained on it) along with four probe tasks that, to our knowledge, neither model
was trained on (HellaSwag, ARC-Challenge, MMLU, TruthfulQA).
Table~\ref{tab:behavior_holdout} reports a significant negative random-effects correlation across all held-out
probes in Llama. Flan HellaSwag is instead significantly positive. Empty, immediate
end-of-sequence (EOS) responses are observed output and are represented as an explicit empty-output
semantic state in the primary analysis. In the Llama full-FT runs,
38.64\% of HellaSwag leaf outputs were empty, versus 0.02\% for TruthfulQA and 0.00\% for ARC-C
and MMLU. For HellaSwag, including immediate EOS gives $r_{\mathrm{DL}}=-0.17$ [$-0.24,-0.10$], whereas excluding empty completions strengthens the all-node estimate to $-0.27$ [$-0.34,-0.19$] and leaves the other probe estimates essentially unchanged (App.~\ref{app:empty-sensitivity}). These results show that shared fine-tuning ancestry is associated with an increase in behavioral similarity on held-out probes. The near-zero Flan held-out-translation cell is instead a degeneration diagnostic along with three other probes for Flan (App.~\ref{app:translation_probe}).
For Flan especially, classification-trained leaves collapse to their label vocabulary after repeated fine-tuning and so we mark them as degenerate outputs.



\begin{table}[t]
\caption{Weight$\to$behavior bridge under full fine-tuning: random-effects pooled
  $r_{\mathrm{DL}}$ with 95\% CI. Immediate EOS is represented explicitly; $^\ddagger$degenerate
  Flan probe. Heterogeneity is in App.~\ref{app:regression}.}
\label{tab:behavior_holdout}
\centering
\small
\setlength{\tabcolsep}{5pt}
\begin{tabular}{lcccccc}
\toprule
$r_{\mathrm{DL}}$ [95\% CI] & \multicolumn{2}{c}{Translation} & HellaSwag & ARC-C & MMLU & TruthfulQA \\
\cmidrule(lr){2-3}
Model & trained & held-out & held-out & held-out & held-out & held-out \\
\midrule
Llama-3.2-1B & \shortstack{$-0.20$\\{\scriptsize$[-.26,-.13]$}} & \shortstack{$-0.23$\\{\scriptsize$[-.29,-.17]$}} & \shortstack{$-0.17$\\{\scriptsize$[-.24,-.10]$}} & \shortstack{$-0.20$\\{\scriptsize$[-.26,-.13]$}} & \shortstack{$-0.19$\\{\scriptsize$[-.26,-.12]$}} & \shortstack{$-0.22$\\{\scriptsize$[-.26,-.16]$}} \\
\addlinespace
Flan-T5-base & \shortstack{$-0.10$\\{\scriptsize$[-.16,-.04]$}} & \shortstack{$+0.04^{\ddagger}$\\{\scriptsize$[-.01,+.09]$}} & \shortstack{$+0.09$\\{\scriptsize$[+.05,+.14]$}} & \shortstack{$-0.14^{\ddagger}$\\{\scriptsize$[-.19,-.09]$}} & \shortstack{$-0.16^{\ddagger}$\\{\scriptsize$[-.20,-.11]$}} & \shortstack{$-0.17^{\ddagger}$\\{\scriptsize$[-.22,-.12]$}} \\
\bottomrule
\end{tabular}
\end{table}


\paragraph{Comparing to behavior-only lineage.}
\label{sec:results:blackbox}

When weights are unavailable or different architectures need to be compared, lineage must be inferred from behavior alone. We compare our weight-based
estimation to PhyloLM~\citep{yax2025phylolm}, scored against the
same documented lineage. On the real Mistral family (Fig.~\ref{fig:real-trees}), PhyloLM and our approach
both recover all three forks. Their informative-quartet additivity differs ($A=0.84$ vs.\ $0.98$, respectively), indicating that the weight-based distance matrix has a stronger additive signal.

The Mistral forks are large RLHF or direct-preference-optimization (DPO) weight changes that move behavior substantially. Our controlled trees are a more precise test. In some of the LoRA and low-update regimes fine-tuning barely shifts behavior, yet weight distance can still estimate phylogenies. Running PhyloLM on the shared trees against the same ground truth, behavior-only recovery falls well below weight-based recovery. Under full fine-tuning PhyloLM recovers $52\%$ of clades versus $100\%$ for weights, under rank-8 LoRA it recovers $35\%$ versus $100\%$, and under rank-64 LoRA it recovers $24\%$ versus $100\%$. This indicates that behavior-based distances can preserve some lineage signal but not all in some settings. Overall, our weight-based estimation matches behavior-based estimation on this real family and substantially outperforms it in the controlled experiments.

\begin{figure}[t]
  \begin{center}
    \includegraphics[width=\linewidth]{\figpath/hf_mistral_hero_3panel.pdf}
  \end{center}
  \caption{\textbf{Documented Mistral lineage:} truth (left), our weight-based estimate (middle),
    and PhyloLM's output-based estimate (right)~\citep{yax2025phylolm}. Both recover all three splits.}
  \label{fig:real-trees}
\end{figure}
% -----------------------------------------------------------------------
\section{Discussion}
\label{sec:discussion}
% -----------------------------------------------------------------------

We show that model family phylogeny can be estimated from weights alone, and that using a small layer subset---about 3\% of Flan's scalar weights in the encoder-key subset---estimates training trees as well as using all analyzed layers. We find that clade recovery works across encoder--decoder and decoder-only models and survives LoRA training at low rank ($r=8$). Across our controlled conditions, PAER ($\mathrm{FN}=0$) ranges from 85--100\%. The label-free additivity proxy $A$ is associated with increased clade recovery in our experiments, although the relationship depends on architecture and layer subset. In addition to estimating phylogeny and direct lineage from weights, we show that models closer in weight space tend to have more similar held-out outputs. The effect depends on the probe and ranges from $r_{\mathrm{DL}}=-0.23$ to $-0.17$ across interpretable held-out Llama probes. Finally, weight- and behavior-based estimation recover the same documented Mistral family, while weight-based estimation is substantially stronger on several controlled trees and requires no prompts.


We note that our experiments and methods have limitations.  One limitation is that our methods for phylogeny estimation using weights currently only work when the underlying architecture is identical, i.e., the weight space is the same.  This means that Llama cannot directly be compared to Mistral, for example. Moreover, our method assumes that all models are leaves unlike \cite{horwitz2024mother}. Another issue is that our controlled training tree generation results from a fairly simple process that may not result in the most realistic training trees.  We attempt to mitigate that with our real-world example. Future directions include better relating models across different architectures such as with CKA and further exploring the mechanism by which lineage is inherited. In summary, we hope that our work illuminates the need to consider models holistically rather than only as a collection of benchmark scores, and shows how analyzing weight space can shed light on model behavior.


% -----------------------------------------------------------------------
% Author acknowledgments omitted for double-blind review; add at camera-ready.

\subsection*{AI use statement}
Generative AI assistants were used to write and debug software across the experimental
pipeline (data processing, weight-distance computation, tree reconstruction and scoring, statistical
analysis, figures, and tables) and to review and edit the manuscript, including checking consistency between numerical results and prose,
improving readability, searching and summarizing relevant literature, and formatting citations.  The research questions, experimental design,
hypotheses, conceptual framework, and scientific interpretations are the authors' own. The authors
reviewed and tested AI-assisted code and independently verified all results, figures, numerical claims,
citations, and manuscript text. The authors take responsibility for the final content of this work,
including all text, claims, and artifacts produced with the aid of generative AI.

% -----------------------------------------------------------------------
\bibliography{references}
\bibliographystyle{iclr2027_conference}


% -----------------------------------------------------------------------
\appendix
% -----------------------------------------------------------------------

\section{Experimental Details}
\label{app:details}

\subsection*{Tree Generation}
\label{app:tree-gen}

Training trees are generated via a Poisson branching process as follows:

\begin{enumerate}
    \item The virtual root draws $k \sim \mathrm{Poisson}(\lambda = 1.5)$ children, with a minimum of 1. We instantiate the same accepted topologies and task/data assignments with either Flan-T5-base or Llama-3.2-1B.
    \item Each queued node below the depth limit draws $k \sim \mathrm{Poisson}(\lambda = 1.5)$ children while node budget remains. Nodes that draw $k = 0$ naturally become leaves; no separate pruning step is applied.
    \item Expansion is breadth-first, with a budget of at most $n_\mathrm{nodes} = 14$ training nodes (not counting the root) and maximum depth $n_\mathrm{layers} = 4$. Reaching depth 4 on one branch prevents only those depth-4 nodes from expanding; shallower nodes already in the queue can still produce children while budget remains. Growth ends when the queue is empty or the node budget is exhausted. Candidate trees with fewer than four leaves are rejected, and seeded draws continue until 50 trees are accepted.
    \item Each non-root node is assigned one of 36 task/dataset pairs (9 summarization, 10 classification, 8 QA, and 9 translation). The pairs are sampled without replacement within a tree, so datasets do not repeat but task families can. Topology generation and task/dataset assignment are seeded separately.
    \item Topology recovery uses only accepted trees with at least one nontrivial unrooted truth split---an internal edge that partitions the leaves into two groups of at least two. Four accepted trees each have four leaves but reduce to unresolved stars, with only trivial one-versus-three splits; this leaves 46 trees for recovery. Four leaves alone is not exclusionary.
\end{enumerate}

This procedure produces trees with 4--10 observed leaves and variable depth and branching
structure. In our recovery experiments, we use the 46 trees that have at least one nontrivial
unrooted truth split. Across these 46 trees, there are 609 trained nodes (349 leaves and 260 internal
nodes), in addition to 46 untrained roots. The ground-truth topology is fully known by construction,
enabling exact RF evaluation.

\paragraph{Experiment structure.} All topology-recovery comparisons use the same 46 retained
unlabeled tree topologies, with node identities aligned across conditions. In the ordinary Flan
and Llama cohorts, every aligned node also retains the same task/dataset label across root models
and fine-tuning interventions. The matched no-translation cohorts preserve those topologies
and node identities, but replace translation assignments with unused classification, QA, or
summarization datasets (App.~\ref{app:translation_probe}). Thus, comparisons are paired by tree and,
except for the intended translation-removal intervention, by node label.

Different analyses use fixed eligible subsets of these shared trees. Topology recovery uses the
common 46 trees with at least one nontrivial unrooted truth split (excluding topology IDs 015, 020,
037, and 047), while direct ordering uses the common 26 trees with a defined cross-root-branch
comparison. PhyloLM produces a behavioral reconstruction for every tree in each matched Llama
condition, and its topology summaries and the matched weight values in
Table~\ref{tab:variants} use the same 46-tree recovery subset. The Flan and Llama per-layer profiles
in Fig.~\ref{fig:layertrace_grid} likewise use their matched 46-tree full-fine-tuning cohorts.

\subsection*{Training Configuration}
Table~\ref{tab:training_config} summarizes the audited configurations for
\texttt{google/flan-t5-base}\footnote{Pinned revision:
\nolinkurl{7bcac572ce56db69c1ea7c8af255c5d7c9672fc2}.} and the base, non-Instruct
\texttt{meta-llama/Llama-3.2-1B}\footnote{Pinned revision:
\nolinkurl{4e20de362430cd3b72f300e6b0f18e50e7166e08}.} Each child is initialized
from its parent and trained under the configuration for its cohort.

\begin{center}
\begin{minipage}{\linewidth}
\centering
\refstepcounter{table}
\label{tab:training_config}
\textbf{Table \thetable:} Training configuration for the controlled Flan and Llama cohorts.
\vspace{4pt}

\small
\setlength{\tabcolsep}{4pt}
\begin{tabular}{@{}p{0.19\linewidth}p{0.38\linewidth}p{0.38\linewidth}@{}}
\toprule
Setting & Flan-T5-base & Llama-3.2-1B \\
\midrule
Objective & Sequence-to-sequence & Completion-only causal language modeling \\
Training length & 2,000 steps; no early stopping & 2,000 steps; no early stopping \\
Learning rate & $3\times10^{-4}$ (full FT and LoRA) & $2\times10^{-5}$ (full FT); $3\times10^{-4}$ (LoRA) \\
Batching & Batch size 8; gradient accumulation 4; effective batch 32 & Batch size 8; gradient accumulation 4; effective batch 32 \\
Schedule & 200 warmup steps; weight decay 0.01 & 200 warmup steps; weight decay 0.01 \\
Precision & bfloat16 & bfloat16 \\
Sequence limits & 512 input tokens; 128 target tokens & 1,024 tokens total for prompt plus completion \\
Padding & Right & Right \\
LoRA & $r=8$, $\alpha=16$, dropout 0.05; module scopes in Table~\ref{tab:lora_variants} & QKV; $(r,\alpha)=(8,16)$ or $(64,128)$; dropout 0.05 \\
Data sample & At most 10,000 train and 1,000 evaluation examples per dataset & At most 10,000 train and 1,000 evaluation examples per dataset \\
\bottomrule
\end{tabular}
\end{minipage}
\end{center}

Training and dataset subsampling are seeded and held fixed across matched conditions and across
nodes assigned the same dataset. Topology generation and task/dataset assignment use separate seed
streams, as described above.

\subsection*{Behavioral Probes and Analysis Population}
Within each behavioral cohort, the primary analysis uses all 641 trained, non-root checkpoints
across 50 trees: 276 trained internal nodes and 365 terminal leaves. Each internal node is an actual
checkpoint saved during training, not a latent ancestor inferred from its descendants; we evaluate
its weights and outputs directly, just as we do for a leaf. The leaf-only analysis is reported
separately as a sensitivity check. For each within-tree checkpoint pair, weight distance is the mean
active-layer cosine distance described above. Translation uses 200 frozen prompts with one
deterministic greedy completion per checkpoint (up to 128 new tokens for Flan and 64 for Llama).
HellaSwag, ARC-Challenge, MMLU, and TruthfulQA
each use 100 frozen prompts and three sampled completions per prompt, with a distinct generation seed
for each draw (temperature 1, $p=1$; up to 64 new tokens). Dolly, used only in the
end-of-sequence (EOS) sensitivity, uses 30 prompts, three draws, and up to 96 new tokens.

Outputs are embedded with \texttt{all-MiniLM-L6-v2}\footnote{Pinned revision:
\nolinkurl{1110a243fdf4706b3f48f1d95db1a4f5529b4d41}.}. We average prompt-paired cosine similarity (and
prompt-by-draw similarity for the three-draw probes) before correlating it with weight distance within
each tree. The inference unit is the tree. Within-tree Fisher-$z$ variances are estimated by deleting one
trained node at a time, after which the tree correlations are pooled by
DerSimonian--Laird~\citep{dersimonianlaird1986}. Pair counts are descriptive and are not treated as
independent observations.

\section{Layer Variant Definitions}
\label{app:layers}
Two families of variants appear in the paper. \emph{layer subsets} (Table~\ref{tab:layer_subsets})
restrict the full-weight distance to a chosen set of layers, testing which layers carry the lineage
signal; they are defined on Flan-T5-base (12 encoder + 12 decoder blocks). \emph{LoRA variants}
(Table~\ref{tab:lora_variants}) instead change which modules are fine-tuned, all at rank $r=8$,
$\alpha=16$, dropout $0.05$. Their controlled-tree analysis merges the full
root-to-node adapter chain so distances are computed between standalone full-weight leaves. The
all-projections row is retained as an explicitly marked cumulative-update ablation: it compares the
cumulative root-to-node adapter displacement with the shared base omitted, whereas the other LoRA rows
use standalone merged full weights. It is therefore excluded from the merged-weight diagnostic fit.

\begin{table}[h]
\caption{layer subsets (matched-topology Flan-T5-base full FT; cosine; $n=46$). ``SA.k'' is the self-attention key projection;
``pre-attention LayerNorm'' is the \texttt{layer.0.layer\_norm} preceding self-attention. The table's
percentages count layers in the 282-layer trained set, not scalar parameters. In particular,
\texttt{enc\_sak} is 12/282 layers (4.3\%) but 7,077,888/247,577,856 scalar weights (2.86\%).
\emph{Clade} and PAER are polytomy-aware recovery percentages with SEs over the same runs.
\texttt{low\_signal} is the predefined LayerNorm control subset.}
\label{tab:layer_subsets}
\centering
\scriptsize
\setlength{\tabcolsep}{2pt}
\begin{tabular}{llrrrr}
\toprule
Variant & layers included & Count & \% layers & Clade (SE) & PAER (SE) \\
\midrule
\texttt{full}         & all trained layers                                   & 282 & 100          & 100.0 (0.0) & 100.0 (0.0) \\
\texttt{high\_signal} & SA.k $+$ pre-attn LN (encoder $+$ decoder)            & 48  & $\sim$17     & 100.0 (0.0) & 100.0 (0.0) \\
\texttt{sak}          & SA.k only (encoder $+$ decoder)                       & 24  & $\sim$8.5    & 100.0 (0.0) & 100.0 (0.0) \\
\texttt{encoder\_h}   & encoder SA.k $+$ pre-attn LN                          & 24  & $\sim$8.5    & 100.0 (0.0) & 100.0 (0.0) \\
\texttt{enc\_sak}     & encoder SA.k only                                     & 12  & 4.3           & 100.0 (0.0) & 100.0 (0.0) \\
\texttt{low\_signal}  & post-attn/cross-attn/post-FFN LNs (control)           & 36 & $\sim$13 & 100.0 (0.0) & 100.0 (0.0) \\
\bottomrule
\end{tabular}
\end{table}

\begin{table}[h]
\caption{LoRA target-module variants in the fixed-2,000-step Flan-T5-base cohort
($r=8$, $\alpha=16$, dropout $0.05$). Modules $q,k,v,o$ are attention
query/key/value/output projections; \texttt{wi\_0}, \texttt{wi\_1}, and \texttt{wo} are the two gated
FFN inputs and the FFN output. Counts follow the modules actually resolved by PEFT (36 modules per
attention projection and 24 per FFN projection); percentages use the 247,577,856 scalar-weight analysis
denominator.}
\label{tab:lora_variants}
\centering
\begin{tabular}{ll r r}
\toprule
Variant & Target modules & Trainable params & \% scalars \\
\midrule
\texttt{lora\_k\_r8}          & $k$                                  & 442,368   & 0.18 \\
\texttt{lora\_qkv\_r8}        & $q,k,v$                              & 1,327,104 & 0.54 \\
\texttt{lora\_qkvo\_r8}       & $q,k,v,o$                            & 1,769,472 & 0.71 \\
\texttt{lora\_all\_proj\_r8}  & $q,k,v,o,\mathrm{wi}_{0,1},\mathrm{wo}$ & 3,391,488 & 1.37 \\
\bottomrule
\end{tabular}
\end{table}

\subsection*{Per-layer Recovery}
\label{app:layertrace}
Treating each layer as an independent estimator, Fig.~\ref{fig:layertrace_grid} reports the per-layer
FN-rate across architectures. On the Flan trees the signal is diffuse across individual layers:
attention key matrices and the remaining layers each miss about 4\% of true clades on average. This
does not conflict with Table~\ref{tab:layer_subsets}, where averaging only the 12 encoder
\textsc{SelfAttention.k} matrices matches the all-weight reconstruction. \textsc{LayerNorm} produces
prominent local maxima in both architectures.

\begin{figure}[h]
  \begin{center}
    \includegraphics[width=\linewidth]{\figpath/fig3_layertrace_grid.pdf}
  \end{center}
  \caption{Block-aligned per-layer FN-rate (\% of true clades missed; mean $\pm$ 2\,SE,
    cosine, full fine-tuning) for the encoder and decoder of Flan-T5-base ($d{=}768$, $n=46$)
    and the decoder of the matched Llama-3.2-1B cohort ($d{=}2048$, $n=46$). Each transformer block
    occupies one x-position; gray points show the remaining layers without connecting unrelated
    layer families. Flan key matrices and other layers each average about 4\% missed clades, whereas
    the Llama profile is more key-concentrated (1.1\% vs.\ 24.0\%). Color and shape highlight
    self-attention keys, cross-attention keys, and \textsc{LayerNorm}. The two Flan block-0
    relative-attention-bias layers (32.8\% encoder; 55.4\% decoder) are omitted from the plot;
    panel plotting windows extend from $-5\%$ to 5\%, 30\%, and 102\%, respectively.}
  \label{fig:layertrace_grid}
\end{figure}

\section{Blackbox Experiment Details}
\label{app:blackbox}

\paragraph{PhyloLM reproduction.} For the behavior-only comparison on the documented Mistral family
(Fig.~\ref{fig:real-trees}), we reproduce PhyloLM \citep{yax2025phylolm} faithfully using its released
\texttt{math} gene bank and its published Nei-similarity genetic distance (upstream commit
\texttt{8c70edf}). We sample $G=128$ genes and, for each, draw $N=32$ raw completions of $4$ tokens per
model (temperature $1$, no chat template, matching PhyloLM's design), take the $4$-character allele,
form the per-gene allele populations, compute the Nei similarity between every pair of the ten
checkpoints, convert to a distance ($-\log$ of the similarity), and build the tree by neighbor-joining.
The reconstruction is scored against the same documented truth tree and with the same polytomy-aware
scorer used for the whitebox families.

The controlled comparison uses the same locked PhyloLM contract ($G=128$, $N=32$, four-token
continuations, temperature $1$) with sampling seed \texttt{20260803} on the 50 matched-topology Llama
rank-8 LoRA, rank-64 LoRA, and full-FT conditions. For LoRA leaves, we reconstruct and query the
ordered cumulative root-to-leaf adapter chain; full-FT leaves are loaded directly. Each condition
contains 50 completed reconstructions and 365 leaf populations, and all generated alleles are
nonempty. The four trees whose truths contain no nontrivial unrooted split are valid behavioral
reconstructions but cannot define clade recovery, PAER, RF, or FN, leaving the same $n=46$ topology
sample in every row of Table~\ref{tab:variants}.

\section{Additional Organic Lineage Checks}
\label{app:organic-additional}

The paper evaluates three documented organic HuggingFace lineages in total. The ten-checkpoint
Mistral-7B family is the main-text example because it contains three documented nontrivial splits and
supports the matched PhyloLM comparison. Two smaller checks establish that its whitebox result is not
unique to Mistral. In a seven-checkpoint Flan-T5-base lineage, averaged-matrix NJ and the BHV Fr\'echet
mean recover the sole documented Flan-Alpaca descendant clade (FN $=0/1$); its complementary split has
98.2\% Fr\'echet support under cosine distance. In a ten-checkpoint Llama-2-7B lineage, the cosine
Fr\'echet estimate likewise recovers the sole documented seven-model chat-descendant clade (FN
$=0/1$), with 95.2\% support for the complementary base/Nous-Hermes/Vicuna split. These two references
are bushy and specify no ordering within the recovered clade, so we claim recovery of one documented
split in each, not complete resolution of their descendant histories.


\section{Behavioral Meta-Analysis Diagnostics}
\label{app:regression}

\paragraph{Pooling estimator.} Because each tree has its own weight-space scale, a raw pooled slope is
not meaningful across trees. We standardize within tree and estimate the variance of each Fisher-$z$
correlation with a delete-one-trained-node jackknife, then pool tree-level effects using a
DerSimonian--Laird random-effects meta-analysis. Pair counts are descriptive and are never treated as
independent observations.

\paragraph{Per-cell heterogeneity.} Table~\ref{tab:behavior_het} gives, for every cell of the main-text
bridge (Table~\ref{tab:behavior_holdout}, paired semantic DV), $r_{\mathrm{DL}}$ with its $95\%$ CI,
the Higgins $I^2$, and the number of trees $k$. Heterogeneity ranges from low to substantial ($I^2$ $0$--$36\%$).
All cells except Flan translation-withheld exclude zero. The Llama cells and three flagged Flan cells
are negative, Flan HellaSwag is significantly positive, and Flan translation-withheld is near zero.
Flagged Flan probes are output-degeneration diagnostics rather than clean bridge evidence.

\begin{table}[h]
\caption{Per-cell heterogeneity for the weight$\to$behavior bridge (Table~\ref{tab:behavior_holdout};
  paired semantic DV with immediate-EOS completions represented explicitly, full fine-tuning).
  DerSimonian--Laird random-effects $r_{\mathrm{DL}}$ with $95\%$ CI, Higgins $I^2$,
  and number of trees $k$. $^\ddagger$degenerate Flan probe (App.~\ref{app:translation_probe}).}
\label{tab:behavior_het}
\begin{center}
\small
\begin{tabular}{llrrr}
\toprule
Model & Probe & $r_{\mathrm{DL}}$ [95\% CI] & $I^2$ & $k$ \\
\midrule
Llama-3.2-1B & Translation (trained)              & $-0.20$ [$-0.26,-0.13$] & 13\% & 50 \\
Llama-3.2-1B & Translation (held-out)             & $-0.23$ [$-0.29,-0.17$] & 0\% & 50 \\
Llama-3.2-1B & HellaSwag (held-out)               & $-0.17$ [$-0.24,-0.10$] & 26\% & 50 \\
Llama-3.2-1B & ARC-C (held-out)                   & $-0.20$ [$-0.26,-0.13$] & 33\% & 50 \\
Llama-3.2-1B & MMLU (held-out)                    & $-0.19$ [$-0.26,-0.12$] & 36\% & 50 \\
Llama-3.2-1B & TruthfulQA (held-out)              & $-0.22$ [$-0.27,-0.16$] & 19\% & 50 \\
Flan-T5-base & Translation (trained)              & $-0.10$ [$-0.16,-0.04$] & 0\% & 50 \\
Flan-T5-base & Translation (withheld)$^\ddagger$ & $+0.04$ [$-0.01,+0.09$] & 0\% & 50 \\
Flan-T5-base & HellaSwag (held-out)               & $+0.09$ [$+0.05,+0.14$] & 0\% & 50 \\
Flan-T5-base & ARC-C (held-out)$^\ddagger$         & $-0.14$ [$-0.19,-0.09$] & 0\% & 50 \\
Flan-T5-base & MMLU (held-out)$^\ddagger$          & $-0.16$ [$-0.20,-0.11$] & 0\% & 50 \\
Flan-T5-base & TruthfulQA (held-out)$^\ddagger$    & $-0.17$ [$-0.22,-0.12$] & 3\% & 50 \\
\bottomrule
\end{tabular}
\end{center}
\end{table}


\section{Immediate End-of-Sequence (EOS) Sensitivity and Depth}
\label{app:empty-sensitivity}

An immediate EOS is an observed model response rather than a missing observation. The primary semantic
analysis therefore represents an empty completion as an explicit semantic state. Table~\ref{tab:eos-sensitivity}
reports the complementary sensitivity analysis that removes pairs involving an empty completion. The choice
matters materially only for HellaSwag, where immediate EOS is common.

\begin{table}[h]
\caption{Llama-3.2-1B full-FT all-trained-node semantic bridge with immediate-EOS outputs included as an explicit
  semantic state (primary) or excluded (sensitivity). Values are $r_{\mathrm{DL}}$ [95\% CI]. Empty
  rates are pooled over all leaf outputs and all three draws.}
\label{tab:eos-sensitivity}
\centering
\small
\begin{tabular}{lrrr}
\toprule
Probe & EOS included & EOS excluded & Empty outputs \\
\midrule
ARC-Challenge & $-0.195$ [$-0.257,-0.131$] & $-0.195$ [$-0.257,-0.131$] & $0.00\%$ \\
Dolly open-ended & $-0.081$ [$-0.134,-0.028$] & $-0.083$ [$-0.136,-0.030$] & $0.16\%$ \\
HellaSwag & $-0.175$ [$-0.245,-0.103$] & $-0.265$ [$-0.338,-0.189$] & $38.64\%$ \\
MMLU & $-0.189$ [$-0.256,-0.121$] & $-0.189$ [$-0.256,-0.121$] & $0.00\%$ \\
TruthfulQA & $-0.218$ [$-0.275,-0.161$] & $-0.219$ [$-0.275,-0.161$] & $0.02\%$ \\
\bottomrule
\end{tabular}
\end{table}

Table~\ref{tab:eos-depth} reports the same empty-output rate by leaf depth. HellaSwag empties increase
with depth, from 16.72\% at depth 1 to 40.25\% at depth 4. Immediate termination is therefore itself
associated with accumulated fine-tuning depth; excluding these outputs nevertheless strengthens the
negative all-node semantic association, as shown in Table~\ref{tab:eos-sensitivity}.

\begin{table}[h]
\caption{Immediate-EOS output rate by leaf depth in the Llama-3.2-1B full-FT evaluation trees. Rates
  pool outputs across all leaves observed at each depth and all three draws.}
\label{tab:eos-depth}
\centering
\small
\begin{tabular}{lrrrr}
\toprule
Probe & Depth 1 & Depth 2 & Depth 3 & Depth 4 \\
\midrule
ARC-Challenge & 0.00\% & 0.00\% & 0.00\% & 0.00\% \\
Dolly open-ended & 1.45\% & 0.34\% & 0.10\% & 0.06\% \\
HellaSwag & 16.72\% & 30.86\% & 41.13\% & 40.25\% \\
MMLU & 0.00\% & 0.00\% & 0.00\% & 0.00\% \\
TruthfulQA & 0.10\% & 0.04\% & 0.01\% & $<0.01\%$ \\
\bottomrule
\end{tabular}
\end{table}


\section{Matched Translation Holdout}
\label{app:translation_probe}

The no-translation cohort is matched to the ordinary 50-tree cohort. It
preserves every topology, node order, and non-translation assignment, replacing only the 162 nodes
originally assigned translation data with unused classification, QA, or summarization datasets. Thus
479 of 641 node assignments are unchanged, no translation assignments remain, and the cohort differs
from the ordinary full-FT cohort only where translation was removed. We use the same 200-prompt
translation endpoint to test whether weight distance predicts behavioral similarity when translation
was never a training task.

The primary all-trained-node estimate is
$r_{\mathrm{DL}}=+0.04$ [$-0.01,+0.09$] across 50 trees, so its interval includes zero. We retain this
cell as a \emph{degeneration diagnostic}, not as evidence that the weight--behavior bridge reverses.
Pairwise semantic similarity can be driven by whether models fail on the held-out task in the same
task-specific way rather than by genealogical proximity. In the matched Llama-3.2-1B
full-fine-tuning cohort, the translation bridge is negative and
strengthens when translation is withheld (Table~\ref{tab:behavior_holdout}). The near-zero Flan-T5
result remains the diagnostic described here rather than a clean reversal of the bridge.


\end{document}
